\section{Methods}\label{sec:methods}

\subsection{Task}\label{sec:task}

The target predicate is
\texttt{behavioral\_stability\_windowed}~\cite{banu2026operon}, whose
verify function
(\texttt{\_verify\_behavioral\_stability\_windowed} in
\texttt{operon\_ai/core/certificate.py}) implements
\[
\mathrm{holds}(\mathbf{x}, \tau) \iff \max_{i} x_i \le \tau,
\]
where $\mathbf{x}$ is a vector of per-window means and $\tau$ is a
stability threshold.  The obligation set is enumerable: every
violating window $(i, x_i)$ is a concrete obligation.

\paragraph{Candidate contract.}  The GEPA candidate is a two-key
dict: an immutable prose scratch-space \texttt{policy\_prompt} and a
mutable numeric knob \texttt{policy\_throttle} $\in [0, 1]$ (the seed
value is the literal string \texttt{"1.0"} --- uniformly failing).
The harness
(\texttt{eval/convergence/synthetic\_signal\_harness.py}) reads the
throttle directly from the dict via \texttt{parse\_throttle}; no
prose parsing is performed and the \texttt{policy\_prompt} component
is never read by the harness.  Two pieces of GEPA configuration
enforce that the experiment exercises exactly one degree of freedom:
(i) a custom \texttt{module\_selector}, \texttt{\_ThrottleOnlySelector},
that always returns \texttt{[policy\_throttle]} so GEPA never
proposes mutations to the prose slot (the default
\texttt{round\_robin} would otherwise alternate between the two
keys, and the adapter would reject any \texttt{policy\_prompt}
mutation because the prose key is not in its \texttt{components}
allowlist); and (ii) a per-component
\texttt{reflection\_prompt\_template},
\texttt{THROTTLE\_REFLECTION\_TEMPLATE}, that instructs the
reflection LM to emit \emph{only} a bare float in $[0.0, 1.0]$
--- no prose, no markdown fences, no units --- because GEPA's
default reflection template wraps the LM's full response in
backticks, which would land as the new candidate value verbatim and
fail to parse as a float.  Both pieces are wired in
\texttt{eval/convergence/theorem\_6\_experiment.py} (selector,
template, and \texttt{MUTABLE\_COMPONENTS}).

The harness then draws
$W \cdot O$ observations from $\mathcal{N}(\text{throttle},
\sigma^2)$ --- $W$ non-overlapping windows of $O$ observations each
--- and returns the per-window means as the theorem's
\texttt{signal\_values}.  The default calibration is
$\tau = 0.5$, $\sigma = 0.08$, $W = 4$, $O = 8$, target throttle
$\approx 0.35$.  Per-rollout RNG is seeded deterministically via
SHA-256 of \texttt{(run\_seed, data\_inst\_id)} so every batch is
reproducible.

\subsection{Arms}\label{sec:arms}

\paragraph{cert-binary (treatment).}
\texttt{OperonCertificateAdapter} bound to
\texttt{behavioral\_stability\_windowed}.  Score
$\in \{0.0, 1.0\}$; reflection feedback is produced by a
task-specific \texttt{stability\_windowed\_obligation\_formatter}
that reads the harness trajectory and emits, in order: a theorem
framing line (\texttt{Theorem: behavioral\_stability\_windowed
[FAILED/HOLDS]}), a conclusion line, per-window means with
\texttt{[VIOLATING]} markers on windows that exceed $\tau$, and a
final \texttt{Unmet obligation} line.  The per-window text is
deliberately content-matched to the active-control arm so that the
only remaining differences between cert-binary and scalar-evidence
are (a) binarization of the score and (b) the theorem framing
line.

\paragraph{scalar-reward (baseline).}
\texttt{ScalarRewardAdapter}.  Score is the graded reward
\[
r(\mathbf{x}, \tau) = \begin{cases}
1.0 & \max \mathbf{x} \le \tau,\\
\tau / \max \mathbf{x} & \text{otherwise},
\end{cases}
\]
which peaks at 1.0 exactly when the theorem holds, ensuring
cross-arm convergence detection uses the same predicate.  Reflection
feedback is the single line \texttt{Score: <value>}.

\paragraph{scalar-with-evidence (active control).}
\texttt{ScalarWithEvidenceAdapter}.  Same graded reward; reflection
feedback adds a list of per-window means annotated with a
\texttt{[VIOLATING]} marker wherever the window exceeds $\tau$,
plus a final line asking the LM to lower every window mean to
$\le \tau$.  Crucially, this arm does \emph{not} name the theorem
(no ``Theorem:'' line), so the only signal it adds over
scalar-reward is evidence content, not theorem framing.  The arm
exists to separate two confounds in a positive cert-binary result:
(a) binarization, and (b) obligation-text content.

\subsection{Convergence criterion}\label{sec:conv}

A run converges at iteration $k$ iff the best candidate's GEPA
validation aggregate score equals 1.0 for three consecutive
iterations starting at $k$.  For all three arms, the best score
is aggregated over a held-out validation set of 16 rollouts.
Because the scalar arms peak at 1.0 only when the theorem holds
(see $r(\cdot)$ above), the threshold is the same across arms.

\subsection{Protocol}\label{sec:protocol}

Each cell is one \emph{(arm, seed)} pair.  Ten seeds per arm, three
arms, $10 \times 3 = 30$ cells.  GEPA is configured with
\texttt{max\_metric\_calls} $= 50 \cdot 16 \cdot 2 = 1600$, batch
size 16 for training and 16 for validation, reflection LM
\texttt{ollama/gemma4} at default settings, and a
\texttt{raise\_on\_exception=False} policy (GEPA failures are
recorded and excluded from statistics with a note).

\subsection{Statistics}\label{sec:stats}

Primary: mutations-to-convergence, reported as the GEPA
iteration number (1-based; GEPA's iteration 0 is the initial
evaluation, not a mutation step) at the first event of the
convergence streak.  Non-converged runs are imputed at
\texttt{budget\_iterations} (50) for the central-tendency
statistics (mean iterations, Mann-Whitney U); the CDF plot
(Section~\ref{sec:results}) is built \emph{without} imputation and
treats non-converged runs as right-censored, so arm curves level
off at the arm's convergence rate rather than at 1.0.
Runs that raised an exception inside GEPA
(\texttt{gepa\_error} non-null) are excluded from all statistics
and reported separately in \texttt{summary.json} under
\texttt{errored\_runs\_excluded}.  Mann-Whitney U
(\texttt{alternative="less"}) compares cert-binary against each
baseline.  Rank-biserial correlation is reported as effect size;
by convention in this paper, positive rank-biserial means the
treatment dominates (fewer iterations).  Convergence rates are
reported with Wilson 95\% confidence intervals, matching the
statistical conventions of
papers~\cite{banu2026benchmarks,banu2026convergence}.

The predeclared success criterion, matching the companion memo's
decision gate, is:
\[
\frac{\bar n_{\text{cert-binary}}}{\min(\bar n_{\text{scalar}},\, \bar n_{\text{scalar-evidence}})} \le 0.75
\quad\text{with}\quad p < 0.05.
\]
Any smaller observed effect is reported as inconclusive rather than
as a win.

\paragraph{Non-determinism.}  The reflection LM is invoked without a
temperature or seed pin, so per-cell convergence is a random
variable and a single sweep yields point estimates, not fixed
values.  All headline claims are therefore checked against an $N=10$
replication of the full protocol
(Section~\ref{sec:variance}); where the single-sweep estimate and
the replication disagree, the replication's distribution is the
reported result.

\subsection{Ablation}\label{sec:ablation}

On seeds 0--4, we rerun cert-binary with a \emph{minimal} obligation
formatter that replaces
\texttt{stability\_windowed\_obligation\_formatter} with a
single-line variant: \texttt{"constraint violated: window $k$"} (no
per-window numbers, no theorem framing).  Because both formatters
run under the cert-binary adapter --- binarization is held constant
--- the comparison isolates \emph{numeric evidence content +
theorem framing within the binary arm}.  If the minimal-formatter
arm matches the default-formatter arm, neither numeric evidence
content nor theorem framing adds signal once a structural
``window $k$ failed'' indicator is present; if it underperforms,
the numeric content or framing is load-bearing.  The
complementary contrast, cert-binary vs.\ scalar-evidence
(Section~\ref{sec:arms}), tests \emph{binarization plus theorem
framing} while holding per-window numeric evidence constant, since
the two arms emit byte-identical evidence text after the first
line by construction.  Read together, the two contrasts decompose
the cert-binary signal into (a) binarization, (b) theorem framing,
and (c) numeric evidence content.
